\documentclass[11pt,a4paper]{article}
\usepackage[margin=1.8cm]{geometry}
\usepackage{amsmath,amssymb,tabularx,array,needspace,fontspec}
\setmainfont{FreeSerif}
\setlength{\parindent}{0pt}
\setlength{\parskip}{0pt}
\renewcommand{\arraystretch}{1.15}
\pagestyle{plain}
% \sq{roman}{question}
\newcommand{\sq}[2]{\Needspace{3\baselineskip}\par\smallskip\noindent\hangindent=7.2em\hangafter=1\makebox[7.2em][l]{\textbf{#1.}}#2\par}

\begin{document}

\begin{center}
  {\Huge\bfseries The Hope Science Academy}\\[3pt]
  {\small Excellence is our standard not goal}\\[2pt]
  {\small 247-E-1, (Opposite LDA School) Johar Town, Lahore. PH: 0300-8194789, 0322-4157001}
\end{center}

\medskip
\noindent
\begin{tabularx}{\linewidth}{|X|X|X|X|}
\hline
\textbf{Student Name:} & \textbf{Roll Num:} & \textbf{Class:} 11th & \textbf{Paper Code:} \\ \hline
\textbf{Subject Name:} Principles of Accounting & \textbf{Time Allowed:} & \textbf{Total Marks:} 1243 & \textbf{Exam Date:} 02-Nov-2026 \\ \hline
\multicolumn{4}{|l|}{\textbf{Exam Syllabus:} FULL BOOK} \\ \hline
\end{tabularx}

\bigskip
\noindent\textbf{Q3. Write short answers of the following questions.}\hfill\textbf{1$\times$433=433}

\sq{I}{Define the term “Accounting”.}
\sq{II}{What is the difference between “Book Keeping” \& “Accounting”?}
\sq{III}{What is meant by the term “Business”.}
\sq{IV}{What is the difference between “Accounting” \& “Accountancy”?}
\sq{V}{Explain the following concepts. A) Debtors /Accounting receivable B) creditors/Accounts payable.}
\sq{VI}{Define the term “Cash Discount”.}
\sq{VII}{What is the difference between “Asset” \& “Liabilities”?}
\sq{VIII}{What is Capital or owner’s equity?}
\sq{IX}{Define the term “Drawing”.}
\sq{X}{Define “Separate entity concept”.}
\sq{XI}{Define “Going concern concept”.}
\sq{XII}{Define “Money measurement concept”.}
\sq{XIII}{Define cost concept in accounting.}
\sq{XIV}{Define “Dual aspect concept”.}
\sq{XV}{Define “Accounting period concept”.}
\sq{XVI}{What do you meant by matching concept of accounting?}
\sq{XVII}{Define “Realization concept”.}
\sq{XVIII}{What do you mean by trade discount?}
\sq{XIX}{Name the three main branches of accounting.}
\sq{XX}{Explain the terms “Sales” \& “Purchases”.}
\sq{XXI}{Explain the terms “Credit Sale” \& “Credit Purchases”.}
\sq{XXII}{What do you understand by the term “Returns outwards” \& “Returns Inwards”?}
\sq{XXIII}{What is the difference between “Cost accounting” \& “Financial accounting”?}
\sq{XXIV}{What do you understand by the terms “Expenses” and “Revenue”.}
\sq{XXV}{Define Book Keeping.}
\sq{XXVI}{Write down the objectives of accounting.}
\sq{XXVII}{Write down four objectives of book keeping.}
\sq{XXVIII}{What does an accountant do in an organization?}
\sq{XXIX}{Define profession.}
\sq{XXX}{Define accountancy.}
\sq{XXXI}{Write down the kinds of business.}
\sq{XXXII}{Define trading business.}
\sq{XXXIII}{Define manufacturing business.}
\sq{XXXIV}{Define service business.}
\sq{XXXV}{Define proprietor.}
\sq{XXXVI}{What is capital?}
\sq{XXXVII}{What are merchandise / goods?}
\sq{XXXVIII}{What is meant by purchases?}
\sq{XXXIX}{What are cash purchases?}
\sq{XL}{Define credit purchases.}
\sq{XLI}{What are cash sales?}
\sq{XLII}{Define credit sale.}
\sq{XLIII}{Define purchases return.}
\sq{XLIV}{Define sales return.}
\sq{XLV}{Define debtor.}
\sq{XLVI}{Define creditors.}
\sq{XLVII}{Define assets.}
\sq{XLVIII}{Define equity.}
\sq{XLIX}{Define liabilities.}
\sq{L}{What is discount?}
\sq{LI}{Define cash discount.}
\sq{LII}{Define trade discount.}
\sq{LIII}{Define depreciation.}
\sq{LIV}{What is meant by accumulated depreciation?}
\sq{LV}{What is meant by expenses?}
\sq{LVI}{What is revenue?}
\sq{LVII}{Define invoice.}
\sq{LVIII}{Define voucher.}
\sq{LIX}{What is meant by cash memo?}
\sq{LX}{Define commission.}
\sq{LXI}{Define preliminary expenses.}
\sq{LXII}{Define cost accounting.}
\sq{LXIII}{What is insolvency?}
\sq{LXIV}{What is cash system of accounting?}
\sq{LXV}{What is accrual system of accounting?}
\sq{LXVI}{What is accounting period concept?}
\sq{LXVII}{Define matching principle.}
\sq{LXVIII}{Define consistency principle.}
\sq{LXIX}{What is meant by “Transaction”?}
\sq{LXX}{Differentiate between “External transaction” \& “Internal transaction”?}
\sq{LXXI}{What is meant by “Paper transaction”?}
\sq{LXXII}{What do you mean by “Monetary Events” \& “Non-Monetary Events”?}
\sq{LXXIII}{Differentiate between “Quantitative changes” \& “Qualitative changes”?}
\sq{LXXIV}{What are the three basic elements of accounting equation?}
\sq{LXXV}{Liabilities of a business Rs:150,000 which are 0.25 of proprietorship what will be the amount of: a)Assets b) Capital.}
\sq{LXXVI}{Write a transaction that will:\\
a. Increase a liability and increase an asset?\\
b. Increase an asset and increase Proprietorship?\\
c. Decrease an asset and increase an asset?\\
d. Decrease an asset and decrease in owner’s equity?}
\sq{LXXVII}{Complete the missing amounts in each of the following cases:\\[2pt]\hspace*{3em}\begin{tabular}{@{}lll@{}}
\multicolumn{3}{@{}l}{Assets = Liabilities + Owner\textquoteright s Equity}\\
9,9500 & ? & 50,500 \\
? & 35,500 & 90,000 \\
75,500 & 40,200 & ? \\
\end{tabular}}
\sq{LXXVIII}{What will be the amount of capital if assets are amounting Rs:50,000 and liabilities Rs:35,000.}
\sq{LXXIX}{Define the term “equity”.}
\sq{LXXX}{What is cash transaction?}
\sq{LXXXI}{What is credit transaction?}
\sq{LXXXII}{What is meant by internal transaction?}
\sq{LXXXIII}{What is meant by accounting equation?}
\sq{LXXXIV}{Write down the basic equation of accounting.}
\sq{LXXXV}{Name the element present on equity side of an accounting equation.}
\sq{LXXXVI}{What is the effect of increase in income on capital?}
\sq{LXXXVII}{If assets are Rs. 80,000 and liabilities are Rs. 50,000, identify the amount of capital.}
\sq{LXXXVIII}{What is the effect of increase in expense on capital?}
\sq{LXXXIX}{If assets are Rs. 60,000 and capital is Rs. 40,000, then what will be the liabilities?}
\sq{XC}{What will be the amount of liabilities if assets are Rs. 45,000 and capital is Rs. 20,000?}
\sq{XCI}{If liabilities are Rs. 30,000 and assets are Rs. 80,000, find capital.}
\sq{XCII}{If assets are Rs. 120,000 and liabilities are Rs. 30,000, identify the amount of capital.}
\sq{XCIII}{If liabilities are Rs. 15,000 and capital is Rs. 35,000, then what will be the assets?}
\sq{XCIV}{Define “Double entry system”.}
\sq{XCV}{What do you mean by Single entry system?}
\sq{XCVI}{Differentiate between “Double entry system” and “Single entry system”.}
\sq{XCVII}{Define the term “Account”.}
\sq{XCVIII}{How many types of Accounts are there?}
\sq{XCIX}{What are “Nominal Accounts”?}
\sq{C}{Define “Real or Property accounts”.}
\sq{CI}{What are “Personal account”?}
\sq{CII}{Discuss the rules of debit and credit for Personal, Real and Nominal accounts?}
\sq{CIII}{Rules of debit \& credit for assets \& expenses?}
\sq{CIV}{Rules of debit \& credit for liabilities, revenue \& capital?}
\sq{CV}{Analyze nominal account.}
\sq{CVI}{“Purchase goods on account” analyze the transaction.}
\sq{CVII}{“Cash received from Nabeel” analyze the transaction.}
\sq{CVIII}{Define the term “Journal”.}
\sq{CIX}{Differentiate between Journal and Journalizing.}
\sq{CX}{What is meant by Entry?}
\sq{CXI}{What do you understand by the term Simple entry?}
\sq{CXII}{What is compound entry?}
\sq{CXIII}{Differentiate between “Simple entry” and “Compound entry”.}
\sq{CXIV}{What is meant by “Narration”?}
\sq{CXV}{Following accounts are being maintained in the books of X. Classify them under Assets, Liabilities, Owner’s equity, Expenses and Revenue headings: Interest (Cr), Bank, Commission, Rent, Discount( Cr), Depreciation, Drawings, A’s loan account (Cr), Furniture \& fitting, Capital, Bad debts, Goodwill, Wages outstanding, stock}
\sq{CXVI}{What is journalizing?}
\sq{CXVII}{Why is journal called the book of original entry?}
\sq{CXVIII}{Give two advantages of journal.}
\sq{CXIX}{What is meant by ledger folio?}
\sq{CXX}{Draw the format of journal.}
\sq{CXXI}{Calculate trade discount at 10\% of Rs. 400.}
\sq{CXXII}{Calculate trade discount at 10\% of Rs. 8,000.}
\sq{CXXIII}{Calculate 10\% trade discount on sales of Rs. 50,000.}
\sq{CXXIV}{Pass the journal entry for goods given as charity Rs. 5,000.}
\sq{CXXV}{Pass any compound entry.}
\sq{CXXVI}{Goods withdrawn by owner for personal use. Pass journal entry.}
\sq{CXXVII}{Pass the journal entry for purchase of building Rs. 8,000.}
\sq{CXXVIII}{Goods purchased for Rs. 500. Give journal entry.}
\sq{CXXIX}{Cash deposited Rs. 350 in bank. Pass journal entry.}
\sq{CXXX}{Wages paid to worker. Pass journal entry.}
\sq{CXXXI}{What is meant by the term “Ledger”?}
\sq{CXXXII}{What is the difference between Journal \& Ledger?}
\sq{CXXXIII}{Define the term “Posting”.}
\sq{CXXXIV}{Differentiate between “Posting” \& “Journalizing”.}
\sq{CXXXV}{What are the rules for balancing an account?}
\sq{CXXXVI}{Explain the terms “Debit balance”, “Credit Balance” \& “Zero Balance”.}
\sq{CXXXVII}{What does Ledger Entry or Ledger Posting mean?}
\sq{CXXXVIII}{Distinguish between Ledger entry and Journal entry.}
\sq{CXXXIX}{How many forms of Ledger accounts are there?}
\sq{CXL}{What are the stages of Accounting Cycle?}
\sq{CXLI}{What is balancing?}
\sq{CXLII}{What is “Trail balance”?}
\sq{CXLIII}{What are the methods of preparing Trial Balance?}
\sq{CXLIV}{Why ledger is called “King of all Books of accounts”?}
\sq{CXLV}{What are the normal balances for Assets, Liabilities, Revenues and expenses and owner’s equity?}
\sq{CXLVI}{Explain the “Debtor’s Ledger and “Creditor’s Ledger”?}
\sq{CXLVII}{Draw the format of Ledger account.}
\sq{CXLVIII}{What is Balance?}
\sq{CXLIX}{Define journal folio.}
\sq{CL}{Draw sketch for accounting cycle.}
\sq{CLI}{Why is the ledger called the king of all books of accounts?}
\sq{CLII}{State the merits of trial balance.}
\sq{CLIII}{Describe the process of balancing of accounts.}
\sq{CLIV}{What is debit balance?}
\sq{CLV}{What is credit balance?}
\sq{CLVI}{What is zero balance?}
\sq{CLVII}{Define the bills of exchange.}
\sq{CLVIII}{What is meant by acceptance?}
\sq{CLIX}{Define the terms “Drawer”, “Drawee” \& “Payee”?}
\sq{CLX}{Differentiate between “Trade bills” and “accommodation bills”?}
\sq{CLXI}{Distinguish between inland and Foreign bills?}
\sq{CLXII}{Distinguish between “Term bills” and “Demand bills”.}
\sq{CLXIII}{What is meant by Tenor and Maturity?}
\sq{CLXIV}{Differentiate between dishonor of bill and renewal of a bill.}
\sq{CLXV}{What does discounting of a bill mean?}
\sq{CLXVI}{How many parties are involved in bill of exchange?}
\sq{CLXVII}{What is the difference between bill sent to bank for collection and discounting of a bill?}
\sq{CLXVIII}{What do you mean by endorsement of bill?}
\sq{CLXIX}{Define the term “Days of Grace”.}
\sq{CLXX}{Define the term Rebate.}
\sq{CLXXI}{What is meant by retiring a bill under rebate?}
\sq{CLXXII}{What is meant by dishonor of bill?}
\sq{CLXXIII}{Explain the term “accommodation bill”.}
\sq{CLXXIV}{Define noting charges.}
\sq{CLXXV}{Define the term “Bill receivable” and “Bill payable”.}
\sq{CLXXVI}{Distinguish between dishonor of a bill and cancellation of a bill.}
\sq{CLXXVII}{What is meant by renewal of a bill?}
\sq{CLXXVIII}{Define Promissory Note.}
\sq{CLXXIX}{What is meant by insolvency of a person?}
\sq{CLXXX}{What is meant by negotiable instruments?}
\sq{CLXXXI}{Define bill of exchange.}
\sq{CLXXXII}{What is meant by acceptance of bill of exchange?}
\sq{CLXXXIII}{Name parties involved in a bill of exchange.}
\sq{CLXXXIV}{Define drawer.}
\sq{CLXXXV}{Define drawee.}
\sq{CLXXXVI}{Define payee.}
\sq{CLXXXVII}{Differentiate between bill receivable and bill payable.}
\sq{CLXXXVIII}{What are grace days?}
\sq{CLXXXIX}{Define discounting of bill of exchange.}
\sq{CXC}{Pass journal entry for discounting of bill of exchange.}
\sq{CXCI}{What is meant by renewal of bill of exchange?}
\sq{CXCII}{Define noting charges.}
\sq{CXCIII}{What do you mean by endorsement of bill?}
\sq{CXCIV}{Write two benefits of bill of exchange.}
\sq{CXCV}{What is meant by retiring of bill?}
\sq{CXCVI}{Define rebate.}
\sq{CXCVII}{Define trade bill of exchange.}
\sq{CXCVIII}{Define sight bill.}
\sq{CXCIX}{Define inland bill of exchange.}
\sq{CC}{Define foreign bill of exchange.}
\sq{CCI}{Differentiate between trade bill of exchange and accommodation bill of exchange.}
\sq{CCII}{Define cheque. List down its parties.}
\sq{CCIII}{Name kinds of cheques.}
\sq{CCIV}{Define bearer cheque.}
\sq{CCV}{Define order cheque.}
\sq{CCVI}{Define crossed cheque.}
\sq{CCVII}{Define blank cheque.}
\sq{CCVIII}{What is meant by honour of cheque?}
\sq{CCIX}{What is meant by dishonoured cheque? Give reasons for dishonour of cheque.}
\sq{CCX}{Differentiate between bearer and crossed cheque.}
\sq{CCXI}{Differentiate between bearer and order cheque.}
\sq{CCXII}{Differentiate between trade bill of exchange and cheque.}
\sq{CCXIII}{Define bank draft.}
\sq{CCXIV}{}
\sq{CCXV}{What do you understand by Subsidiary books?}
\sq{CCXVI}{Write down the names of four Subsidiary books.}
\sq{CCXVII}{Define “Cash Book”, or “Cash Journal”.}
\sq{CCXVIII}{What do you understand by Purchases book/Purchases journal?}
\sq{CCXIX}{Explain Sales Book/ Sales Journal.}
\sq{CCXX}{What is a Sales returns journal or Sales returns Book?}
\sq{CCXXI}{What is purchases returns journal or purchases returns Book?}
\sq{CCXXII}{What is bills receivable book?}
\sq{CCXXIII}{What is bills payable book?}
\sq{CCXXIV}{What do you understand by proper journal/ principal journal?}
\sq{CCXXV}{How many groups a journal is mainly divided?}
\sq{CCXXVI}{What is meant by “Special journal”?}
\sq{CCXXVII}{Distinguish between “Sales account” \& “Sales journal”?}
\sq{CCXXVIII}{What is meant by “Debit note”?}
\sq{CCXXIX}{Explain the term “Credit note”.}
\sq{CCXXX}{Differentiate between “Purchases account “and “Purchases journal”.}
\sq{CCXXXI}{What is general journal?}
\sq{CCXXXII}{What is petty cash book?}
\sq{CCXXXIII}{Define the term “Bank”.}
\sq{CCXXXIV}{Write three types of Bank accounts.}
\sq{CCXXXV}{Differentiate between “Current account” \& “Fixed account”.}
\sq{CCXXXVI}{What is saving bank account?}
\sq{CCXXXVII}{What do you understand by the term bank overdraft?}
\sq{CCXXXVIII}{What is a cash book? How does it differ from a cash account?}
\sq{CCXXXIX}{Define “Voucher”.}
\sq{CCXL}{How many books does a banker give to the customer at the time of first deposit?}
\sq{CCXLI}{What do you understand by bank pass book?}
\sq{CCXLII}{Distinguish between Pay-in-slip and cheque book?}
\sq{CCXLIII}{Explain “Endorsement of the cheque”.}
\sq{CCXLIV}{What is meant by dishonor of cheque?}
\sq{CCXLV}{How many classifications of cash book are there?}
\sq{CCXLVI}{Define single column cash book.}
\sq{CCXLVII}{Differentiate between “Double column book” \& “Treble column cash book”.}
\sq{CCXLVIII}{Define the contra entry.}
\sq{CCXLIX}{What is a “Bank Reconciliation statement”?}
\sq{CCL}{Write any three causes of disagreement between the balance as per cash book and pass book.}
\sq{CCLI}{What is a “petty cash book”?}
\sq{CCLII}{What is an imprest system?}
\sq{CCLIII}{Define “uncredited Cheques”.}
\sq{CCLIV}{Define unpresented cheque.}
\sq{CCLV}{What will be the favorable balance and unfavorable balance in each book \& pass book?}
\sq{CCLVI}{Which discount (cash discount, Trade discount) is recorded in the cash book?}
\sq{CCLVII}{What are two major functions of bank?}
\sq{CCLVIII}{Define bank account.}
\sq{CCLIX}{What is current account?}
\sq{CCLX}{Define fixed deposit account.}
\sq{CCLXI}{Which books are issued by the bank to its customers at the time of opening bank account?}
\sq{CCLXII}{Define cheque book.}
\sq{CCLXIII}{Define bank charges.}
\sq{CCLXIV}{Name the kinds of cash book.}
\sq{CCLXV}{Define pay-in slip.}
\sq{CCLXVI}{Define double column cash book.}
\sq{CCLXVII}{What is meant by three column cash book?}
\sq{CCLXVIII}{Cash book has always debit balance, discuss.}
\sq{CCLXIX}{Draw the sketch of three column cash book.}
\sq{CCLXX}{What are petty expenses?}
\sq{CCLXXI}{Define petty cash.}
\sq{CCLXXII}{How can you compare cash book and pass book?}
\sq{CCLXXIII}{Explain two causes of disagreement of pass book and cash book.}
\sq{CCLXXIV}{Define bank statement.}
\sq{CCLXXV}{What are omitted cheques?}
\sq{CCLXXVI}{What is meant by blank cheque?}
\sq{CCLXXVII}{Define cross cheque.}
\sq{CCLXXVIII}{What is ATM card?}
\sq{CCLXXIX}{Define closing entries.}
\sq{CCLXXX}{Define adjusting entries.}
\sq{CCLXXXI}{Show the complete “Accounting cycle”.}
\sq{CCLXXXII}{Differentiate between “Trading account and “Profit \& loss account”.}
\sq{CCLXXXIII}{Differentiate between “Gross profit” and “Net profit”.}
\sq{CCLXXXIV}{Differentiate between fixes assets and fixed liabilities.}
\sq{CCLXXXV}{Define “Balance Sheet”.}
\sq{CCLXXXVI}{What do you mean by “Real assets”?}
\sq{CCLXXXVII}{How many types of Real assets are there?}
\sq{CCLXXXVIII}{Which assets are called “Intangible assets”?}
\sq{CCLXXXIX}{What do you understand by Current assets, give three examples?}
\sq{CCXC}{What do you understand by current liabilities, give three examples?}
\sq{CCXCI}{How many types of current liabilities are there?}
\sq{CCXCII}{Define “Contingent liabilities”.}
\sq{CCXCIII}{Differentiate between “Intangible assets” and “Fictitious assets”?}
\sq{CCXCIV}{Define “deferred liabilities”.}
\sq{CCXCV}{Define liquid or quick assets.}
\sq{CCXCVI}{Define Liquid or Quick liabilities?}
\sq{CCXCVII}{What is the difference between “Trial Balance” and “Balance sheet”?}
\sq{CCXCVIII}{Define “Wasting assets”.}
\sq{CCXCIX}{Differentiate between “Internal liabilities” and “External Liabilities”.}
\sq{CCC}{Define the term “Capital”.}
\sq{CCCI}{Distinguish between “Direct expenses” and “Indirect expenses” with examples.}
\sq{CCCII}{Distinguish between “Direct revenue” and “Indirect revenue” with examples.}
\sq{CCCIII}{What is the amount of gross profit, if a person buys Rs80,000 worth of goods and sells them for Rs: 1,00,000.}
\sq{CCCIV}{Define final account.}
\sq{CCCV}{Define trading account.}
\sq{CCCVI}{Define direct expenses.}
\sq{CCCVII}{Define gross profit.}
\sq{CCCVIII}{Define gross loss.}
\sq{CCCIX}{Define opening stock.}
\sq{CCCX}{Define closing stock.}
\sq{CCCXI}{Define profit and loss accounts.}
\sq{CCCXII}{Define net profit.}
\sq{CCCXIII}{Define net loss.}
\sq{CCCXIV}{Define indirect expenses with examples.}
\sq{CCCXV}{Define indirect income.}
\sq{CCCXVI}{Define marshalling.}
\sq{CCCXVII}{Define tangible assets.}
\sq{CCCXVIII}{Define fixed assets with examples.}
\sq{CCCXIX}{Define current assets.}
\sq{CCCXX}{Define contingent assets.}
\sq{CCCXXI}{Define fictitious assets.}
\sq{CCCXXII}{Define liabilities.}
\sq{CCCXXIII}{Give two examples of current liabilities.}
\sq{CCCXXIV}{Define long term liabilities.}
\sq{CCCXXV}{Define income statement.}
\sq{CCCXXVI}{Define the term goodwill.}
\sq{CCCXXVII}{What is meant by an adjustment?}
\sq{CCCXXVIII}{What does adjusting entries mean?}
\sq{CCCXXIX}{Define the cash system of accounting.}
\sq{CCCXXX}{Define the accrual or mercantile system of accounting.}
\sq{CCCXXXI}{What do you understand by the term outstanding expenses?}
\sq{CCCXXXII}{What are the prepaid expenses?}
\sq{CCCXXXIII}{What is accrued revenue?}
\sq{CCCXXXIV}{What adjusting entry will be passed for “accrued income”?}
\sq{CCCXXXV}{Pass the adjusting entry for prepaid expenses?}
\sq{CCCXXXVI}{What do you mean by “Unearned revenue”?}
\sq{CCCXXXVII}{Write down the three examples of “Outstanding expenses”.}
\sq{CCCXXXVIII}{Define the term “Depreciation”.}
\sq{CCCXXXIX}{What is meant by accumulated depreciation?}
\sq{CCCXL}{What will be adjusting entry for “Interest on capital”?}
\sq{CCCXLI}{Define the term “debt”.}
\sq{CCCXLII}{What adjusting entry will be for “Interest on drawings”?}
\sq{CCCXLIII}{Pass the adjusting entry for “Depreciation”.}
\sq{CCCXLIV}{What is meant by provision for discount on debtors?}
\sq{CCCXLV}{What will be the adjusting entry for provision for discount on debtors?}
\sq{CCCXLVI}{What do you understand by “Provision for discount on creditors”?}
\sq{CCCXLVII}{What will be the adjusting entry for provision for discount on creditors?}
\sq{CCCXLVIII}{Differentiate between “Bad debts” \& “Doubtful debts”?}
\sq{CCCXLIX}{Pass the adjusting entry for “Bad debts”.}
\sq{CCCL}{Explain the term “Allowance for Uncollectible”.}
\sq{CCCLI}{What is meant by Bad debts recovered?}
\sq{CCCLII}{What will be the adjusting entry for provision for doubtful debts?}
\sq{CCCLIII}{Differentiate between “Normal loss” \& “Abnormal loss”.}
\sq{CCCLIV}{On 1st January the provision for bad debts was Rs: 2000 the bad debts during the year amounted to Rs.950, the debtors on 31ST December were Rs: 40,000, provision for bad debts is to be kept at 5\% of the debtors. Show the effect of given adjustment on profit or loss account.}
\sq{CCCLV}{A employee a team of eight workers who were paid Rs. 3,000 per month each in the year ending 31stDec 2014, At the start of 2015, he paid salaries by 10\% to Rs.3300 per month to each. Show the amount of salaries charged to profit or loss account for the year ending 31stDec 2015.}
\sq{CCCLVI}{Show how the following items will appear in profit and loss account: a) Bad debts RS. 7200 b) Provision for doubtful debts(old) RS 12,000 c) Current year’s provision(new) RS. 3,000.}
\sq{CCCLVII}{Calculate the value of stock in the beginning: sales Rs. 240,000; percentage of gross profit on sales 20\%; purchases Rs. 175,000; closing stock Rs: 30,000.}
\sq{CCCLVIII}{Calculate the value of stock in the beginning: sales Rs: 150,000; purchases Rs. 150,000; closing stock Rs. 60,000; percentage of gross loss on cost 20\%.}
\sq{CCCLIX}{Calculate the value of stock at the end: sales Rs. 450,000; cash purchases Rs. 100,000; credit purchases Rs. 250,000; carriage in Rs. 12,500; gross profit on sales 20\%; opening stock Rs.37,500.}
\sq{CCCLX}{Define two systems of accounting.}
\sq{CCCLXI}{Define accrued income.}
\sq{CCCLXII}{Define unearned income.}
\sq{CCCLXIII}{Define interest on capital.}
\sq{CCCLXIV}{Define bad debts.}
\sq{CCCLXV}{Cost of an asset is Rs. 10,000. Rate of depreciation is 10\% P.A. Calculate depreciation.}
\sq{CCCLXVI}{Cost of an asset is Rs. 15,000. Rate of depreciation is 10\% P.A. Calculate depreciation.}
\sq{CCCLXVII}{What adjusting entry will be passed for outstanding wages?}
\sq{CCCLXVIII}{What adjusting entry will be passed for outstanding salaries?}
\sq{CCCLXIX}{Salary paid in advance Rs. 5,000, pass journal entry.}
\sq{CCCLXX}{What adjusting entry will be passed for prepaid insurance?}
\sq{CCCLXXI}{Pass adjusting entry for depreciation on furniture.}
\sq{CCCLXXII}{What do you mean by work sheet?}
\sq{CCCLXXIII}{Name the five sets of money columns of worksheet.}
\sq{CCCLXXIV}{What are the advantages of worksheet? (any two)}
\sq{CCCLXXV}{Define income statement.}
\sq{CCCLXXVI}{Why balance sheet is to be prepaid?}
\sq{CCCLXXVII}{What is meant by income statement?}
\sq{CCCLXXVIII}{What is the difference between selling expenses and administrative expenses?}
\sq{CCCLXXIX}{What do you understand by “Operating expenses”?}
\sq{CCCLXXX}{How many types of operating expenses are there?}
\sq{CCCLXXXI}{What is a balance sheet?}
\sq{CCCLXXXII}{What do you mean by financial statement?}
\sq{CCCLXXXIII}{What do you mean by cost of goods sold?}
\sq{CCCLXXXIV}{Differentiate between “Recurring expenses” and “Non recurring expenses”?}
\sq{CCCLXXXV}{Explain the term “outstanding expenses”.}
\sq{CCCLXXXVI}{What do you understand by the term “Prepaid Expenses”?}
\sq{CCCLXXXVII}{Define “Depreciation”.}
\sq{CCCLXXXVIII}{What is meant by allowance for doubtful debts?}
\sq{CCCLXXXIX}{Differentiate between “accrued income” and “unearned income”.}
\sq{CCCXC}{Write down the three examples of administrative expenses?}
\sq{CCCXCI}{What do you mean by capital expenditure?}
\sq{CCCXCII}{Describe the revenue expenditures?}
\sq{CCCXCIII}{Give three examples of each: a) Capital expenditures b) Revenue expenditures}
\sq{CCCXCIV}{Differentiate between capital expenditure and revenue expenditures.}
\sq{CCCXCV}{Define deferred revenue expenditures?}
\sq{CCCXCVI}{What do you mean by capitalized expenditures?}
\sq{CCCXCVII}{Differentiate between deferred revenue expenditure and capitalized expenditures.}
\sq{CCCXCVIII}{What are revenue receipts, give two examples?}
\sq{CCCXCIX}{What do you mean by capital receipts give two examples?}
\sq{CD}{Distinguish between capital receipts and revenue receipts.}
\sq{CDI}{What do you mean by capital profit?}
\sq{CDII}{Write down the accounting treatment of capital and revenue profit?}
\sq{CDIII}{What is meant by capital loss, give an example?}
\sq{CDIV}{Distinguish between capital payments and revenue payments?}
\sq{CDV}{Define capital receipts.}
\sq{CDVI}{Define revenue receipt.}
\sq{CDVII}{Give two examples of capital receipts.}
\sq{CDVIII}{Give two examples of revenue receipts.}
\sq{CDIX}{Define capital payment.}
\sq{CDX}{What is meant by revenue payment?}
\sq{CDXI}{Give two examples of capital profit.}
\sq{CDXII}{Define revenue loss.}
\sq{CDXIII}{Differentiate between capital and revenue loss.}
\sq{CDXIV}{What is meant by deferred revenue expenditures?}
\sq{CDXV}{Give four examples of revenue expenditures.}
\sq{CDXVI}{Define error of omission and error of commission.}
\sq{CDXVII}{What is meant by compensating errors?}
\sq{CDXVIII}{What is a suspense account?}
\sq{CDXIX}{Differentiate error of posting and error of casting.}
\sq{CDXX}{If the goods purchased from Ali for Rs.4,000 were recorded in purchase book as Rs.400 and posted to Ali account as Rs.40, what rectifying entry will be passed?}
\sq{CDXXI}{Purchase returns of Rs.250 wrongly posted to sale returns account. How does this affect the trial balance total?}
\sq{CDXXII}{Sales to Ahmad Rs.336 posted to his account as Rs.363, what rectifying entry is to be passed?}
\sq{CDXXIII}{The payment of wages for RS.750 was entered twice in the books by mistake. Pass the rectifying entry.}
\sq{CDXXIV}{Differentiate errors of book keeping and errors of the trial balance.}
\sq{CDXXV}{Pass the rectifying entry, if salary paid for Rs.20,000 has been wrongly debited to rent account.}
\sq{CDXXVI}{Write four kinds of book-keeping errors.}
\sq{CDXXVII}{What is meant by error of posting?}
\sq{CDXXVIII}{Define error of principle.}
\sq{CDXXIX}{Explain dual-sided error.}
\sq{CDXXX}{What are errors of trial balance?}
\sq{CDXXXI}{What are errors of casting?}
\sq{CDXXXII}{Give two examples of errors of principle.}
\sq{CDXXXIII}{Goods sold but omitted to record in books Rs. 3,000. Pass rectifying entry.}

\end{document}
